\documentclass{article}
\usepackage{amsmath, amssymb}
\usepackage[ruled,vlined,linesnumbered]{algorithm2e}

\title{CNN-Adaptive}
\author{-}


\begin{document}

\maketitle

\section{Method 1: Clothoid Path}

\subsection{Path Generation}
The reference path $P(X,Y)=0$ istransitions from a straight line to a 
circle of radius $R$. 
Let $s$ denote the arc length and $L$ the transition length.

Clothoid Spiral ($0 \le s \le L$):  Using the standard Fresnel Formulation:
\[
a = \sqrt{R L \pi},
\]
the path is defined as:
\begin{align}
X(t) &= a \int_{0}^{t} \cos\left(\frac{\pi u^2}{2}\right)\, du \\
Y(t) &= a \int_{0}^{t} \sin\left(\frac{\pi u^2}{2}\right)\, du
\end{align}




\subsection{Optimal Point Identification}

The optimal point $(X^*, Y^*)$ is defined as the orthogonal projection of the 
vehicle position $(X_v, Y_v)$ onto the path.

The optimal parameter $t^*$ is obtained by solving:
\begin{align}
t^* = \arg \min_{t} \; (X_v - X(t))^2 + (Y_v - Y(t))^2
\end{align}
The first order optimality condition is:
\begin{align}
\frac{d}{dt} \left((X_v - X(t))^2 + (Y_v - Y(t))^2\right) \Big|_{t=t^*} = 0
\end{align}
Which simplifies to:
\begin{align}(X_v - X(t^*)) \dot X(t^*) + (Y_v - Y(t^*)) \dot Y(t^*) = 0
\end{align}
And using equations (1) and (2):
\begin{align}(X_v - X(t^*)) \cos\left(\frac{\pi (t^*)^2}{2}\right) + (Y_v - Y(t^*)) 
    \sin\left(\frac{\pi (t^*)^2}{2}\right) = 0
\end{align} 

The solution $t^*$ is obtained numerically (e.g., Newton–Raphson), and:
\[
X^* = X(t^*), \quad Y^* = Y(t^*)
\]


\subsection{Error Extraction}
At $(X^*, Y^*)$, the reference states are:
\begin{align}
\psi_{\mathrm{des}} &= \frac{\pi (t^*)^2}{2} \\
\end{align}

\textbf{Lateral Error:}
\begin{equation}
    D= \sqrt{(X_v - X^*)^2 + (Y_v - Y^*)^2}
\end{equation}
\begin{equation}    
    e_y = \operatorname{sgn} \left( (X_v - X^*)P_Y - (Y_v - Y^*)P_X \right) \cdot D
\end{equation}


\textbf{Heading Error:}
\begin{align}
e_\psi = \operatorname{normalize}(\psi_v - \psi_{\mathrm{des}})
\end{align}






\section{Method 2: Simple Polynomial Path (Cubic Approximation)}

\subsection{Path Generation}
For small deflection angles, the linear increase in curvature can be modeled using a cubic polynomial:
\begin{equation}
Y(X) = k X^3
\end{equation}
where $k$ is a scaling constant. The first and second derivatives are:
\begin{align}
Y'(X) &= 3k X^2 \\
Y''(X) &= 6k X
\end{align}

\subsection{Curvature Analysis}
The exact curvature formula is:
\begin{equation}
\kappa(X) = \frac{Y''(X)}{\left(1 + (Y'(X))^2\right)^{3/2}} = \frac{6kX}{\left(1 + 9k^2 X^4\right)^{3/2}}
\end{equation}
For $X \approx 0$ (the start of the transition), the denominator is approximately $1$, yielding:
\begin{equation}
\kappa(X) \approx 6kX
\end{equation}
This confirms that the curvature increases linearly with the coordinate $X$.

\subsection{Optimal Point Identification}
The optimal point $(X^*, Y^*)$ for a vehicle at $(X_v, Y_v)$ is found by minimizing:
\begin{equation}
X^* = \arg \min_{X} \; (X - X_v)^2 + (k X^3 - Y_v)^2
\end{equation}
The optimality condition (where the distance vector is orthogonal to the tangent) results in a 5th-degree polynomial:
\begin{equation}
(X^* - X_v) + 3k(X^*)^2(k(X^*)^3 - Y_v) = 0
\end{equation}

\subsection{Error Extraction}
The desired heading $\psi_{\mathrm{des}}$ at the optimal point is:
\begin{equation}
\psi_{\mathrm{des}} = \operatorname{atan}(3k (X^*)^2)
\end{equation}
The lateral error $e_y$ and heading error $e_\psi$ are calculated using the same geometric logic as in Method 1.
\end{document}